% 2011 年湖北卷理科第 15 题：n=1,2,3,4 时黑色正方形互不相邻的着色方案。
% 来源 PDF：普通高考/2011/2011湖北理.pdf，SHA-256 ba721f03ae90f88887c5cde27ae98dbea165d33c73e5ba45b4e4f56b0638b075。
% 原页：第 1 页第 2 栏下部，书页码 667。黑格用实心表示。
\begin{tikzpicture}[font=\small, scale=0.42]
  \newcommand{\row}[3]{% #1=n, #2=y, #3=pattern string
    \node[left] at (-0.3,#2) {$n=#1$};
    \foreach \c [count=\i from 0] in {#3}{
      \ifnum\c=1
        \fill (\i,#2) rectangle ++(1,1);
      \else
        \draw (\i,#2) rectangle ++(1,1);
      \fi
    }
  }
  \foreach \pat [count=\j from 0] in {0,1}{
    \begin{scope}[xshift=\j*2.0cm]
      \foreach \c [count=\i from 0] in \pat {
        \ifnum\c=1 \fill (\i,0) rectangle ++(1,1); \else \draw (\i,0) rectangle ++(1,1); \fi
      }
    \end{scope}
  }
  \node[left] at (-0.3,0) {$n=1$};
  % n=2
  \begin{scope}[yshift=-2.0cm]
    \foreach \pat [count=\j from 0] in {00,10,01}{
      \begin{scope}[xshift=\j*2.6cm]
        \foreach \c [count=\i from 0] in \pat {
          \ifnum\c=1 \fill (\i,0) rectangle ++(1,1); \else \draw (\i,0) rectangle ++(1,1); \fi
        }
      \end{scope}
    }
    \node[left] at (-0.3,0) {$n=2$};
  \end{scope}
  % n=3
  \begin{scope}[yshift=-4.0cm]
    \foreach \pat [count=\j from 0] in {000,100,010,001,101}{
      \begin{scope}[xshift=\j*3.4cm]
        \foreach \c [count=\i from 0] in \pat {
          \ifnum\c=1 \fill (\i,0) rectangle ++(1,1); \else \draw (\i,0) rectangle ++(1,1); \fi
        }
      \end{scope}
    }
    \node[left] at (-0.3,0) {$n=3$};
  \end{scope}
  % n=4
  \begin{scope}[yshift=-6.0cm]
    \foreach \pat [count=\j from 0] in {0000,1000,0100,0010,0001,1010,1001,0101}{
      \begin{scope}[xshift=\j*4.6cm]
        \foreach \c [count=\i from 0] in \pat {
          \ifnum\c=1 \fill (\i,0) rectangle ++(1,1); \else \draw (\i,0) rectangle ++(1,1); \fi
        }
      \end{scope}
    }
    \node[left] at (-0.3,0) {$n=4$};
  \end{scope}
\end{tikzpicture}
